\documentclass[11pt]{article}
\usepackage[margin=1in]{geometry}
\usepackage{amsmath,amssymb,amsthm,mathtools}
\usepackage{graphicx}
\usepackage[hidelinks]{hyperref}
\usepackage{xcolor}

\newtheorem{theorem}{Theorem}
\newtheorem{remark}[theorem]{Remark}
\newcommand{\R}{\mathbb R}
\newcommand{\row}{\operatorname{row}}
\setlength{\emergencystretch}{2em}

\title{Partial random circulant kernels are not orthant symmetric\\
\large A probability-one Pluecker obstruction at $M=N/2$}
\author{Autonomous mathematical discovery run \texttt{fa\_banach\_001}}
\date{17 August 2026}

\begin{document}
\maketitle

\noindent\textbf{Status.}
Candidate full negative answer, likely valid pending expert review.

\section{Source question}

Keiper studies partial random circulant measurement matrices generated by a
centered sub-Gaussian vector.  In Section 4 of \cite{Keiper}, PDF pages
13--14, she asks whether the kernel of such a matrix is orthant symmetric.  In
the paper's formulation, if the rows of a random matrix $B$ span the kernel,
one would require
\[
 \mathbb P(BS\in\Omega)=\mathbb P(B\in\Omega)
\]
for every diagonal sign matrix $S$ and measurable set $\Omega$.  In
particular, this condition forces the distribution of the random kernel
subspace itself to be invariant under $K\mapsto KS$.

\begin{center}
\includegraphics[width=0.96\textwidth]{figures/open_problem_crop.png}
\end{center}

The source suspects that the answer is negative and specifically connects the
question to the balanced recovery threshold $M\lesssim N/2$.  The example
below gives a negative answer already at $M=N/2$.

\section{Proof intuition}

The rows of a $2\times4$ partial circulant matrix depend on four random
parameters, but their two-plane is not a generic point of the Grassmannian.
Its six Pluecker coordinates always satisfy an additional quadratic equation
$Q=0$.  This equation is intrinsic to the plane because it is homogeneous in
the Pluecker coordinates.

A single coordinate sign flip changes three of those six coordinates.  The
same quadratic then becomes a nonzero polynomial in the four generating
variables.  With a Gaussian input, a nonzero polynomial vanishes with
probability zero.  Thus an algebraic event has probability one for the
original row space and zero for its sign flip.  Orthogonal complementation
turns the row-space obstruction into the desired kernel obstruction.

\section{Counterexample}

For a two-plane $W\subset\R^4$ with ordered basis $u,v$, write
\[
 p_{ij}(W)=u_i v_j-u_jv_i\qquad(0\leq i<j\leq3)
\]
for its Pluecker coordinates.  A different ordered basis multiplies all six
coordinates by the same nonzero scalar.  Consequently, the zero set of any
homogeneous polynomial in the $p_{ij}$ is a well-defined event on the
Grassmannian.

\begin{theorem}\label{thm:main}
Let $a,b,c,d$ be independent standard real Gaussian random variables and let
\[
 \Phi=\begin{pmatrix}
 a&b&c&d\\
 d&a&b&c
 \end{pmatrix},
\]
the partial circulant matrix formed from rows $0$ and $1$ of the circulant
matrix generated by $(a,b,c,d)$.  Then $\Phi$ has rank two almost surely, and
the law of $\ker\Phi$ is not invariant under coordinate sign changes.  More
strongly, for $S=\operatorname{diag}(-1,1,1,1)$, the laws of
$\row(\Phi)$ and $\row(\Phi)S$ are separated by a measurable algebraic event
having probabilities one and zero.
\end{theorem}

\begin{proof}
The two rows fail to be independent only on the common zero set of the six
$2\times2$ minors.  This set has Gaussian probability zero, since, for
example, the minor $a^2-bd$ is a nonzero polynomial.  Hence the row space is
a two-plane almost surely.

Its six Pluecker coordinates are
\begin{align*}
 p_{01}&=a^2-bd,& p_{02}&=ab-cd,& p_{03}&=ac-d^2,\\
 p_{12}&=b^2-ac,& p_{13}&=bc-ad,& p_{23}&=c^2-bd.
\end{align*}
Define the homogeneous quadratic
\begin{equation}\label{eq:Q}
 Q(p)=(p_{01}-p_{23})(p_{03}+p_{12})-p_{02}^2+p_{13}^2.
\end{equation}
For the coordinates above,
\[
 p_{01}-p_{23}=a^2-c^2,
 \qquad p_{03}+p_{12}=b^2-d^2.
\]
Moreover,
\begin{align*}
 p_{02}^2-p_{13}^2
 &=(ab-cd-bc+ad)(ab-cd+bc-ad)\\
 &=(a-c)(b+d)(a+c)(b-d)\\
 &=(a^2-c^2)(b^2-d^2).
\end{align*}
Therefore
\begin{equation}\label{eq:originalzero}
 Q\bigl(p(\row(\Phi))\bigr)=0
\end{equation}
identically.  Let $\mathcal A$ be the algebraic subset of the Grassmannian
defined by $Q=0$.  Equation \eqref{eq:originalzero} says
$\mathbb P\{\row(\Phi)\in\mathcal A\}=1$.

Now right multiplication by $S=\operatorname{diag}(-1,1,1,1)$ changes the
sign of $p_{01},p_{02},p_{03}$ and leaves $p_{12},p_{13},p_{23}$ unchanged.
Thus
\begin{equation}\label{eq:flippedpoly}
 Q\bigl(p(\row(\Phi)S)\bigr)
 =Q(-p_{01},-p_{02},-p_{03},p_{12},p_{13},p_{23})
\end{equation}
is a polynomial in $a,b,c,d$.  It is not the zero polynomial: at
$(a,b,c,d)=(1,0,1,0)$ its value is $4$.  The zero set of a nonzero real
polynomial has Lebesgue measure zero, and the Gaussian law has a density.
Consequently,
\[
 \mathbb P\{\row(\Phi)S\in\mathcal A\}=0.
\]
The row-space law is therefore not sign-invariant.

Finally, $S$ is orthogonal and
\[
 (\ker\Phi\,S)^\perp=(\ker\Phi)^\perp S=\row(\Phi)S.
\]
If the kernel law were invariant under $K\mapsto KS$, then its orthogonal
complement, the row-space law, would also be invariant.  The preceding
probability-one/probability-zero separation contradicts this.  Hence the
kernel is not orthant symmetric.
\end{proof}

\begin{remark}[Distributional scope]\label{rem:scope}
The proof only uses that $(a,b,c,d)$ has an absolutely continuous joint law;
independence and Gaussianity are a natural match to the source but are not
needed for the algebraic separator.  Thus the same conclusion holds for every
absolutely continuous generating-vector distribution on $\R^4$.
\end{remark}

\section{Independent small-dimensional check}

There is also a direct $2\times3$ witness.  For iid standard Gaussians
$b_0,b_1,b_2$, the kernel of
\[
 \begin{pmatrix}b_0&b_1&b_2\\b_2&b_0&b_1\end{pmatrix}
\]
is spanned almost surely by
\[
 v=(b_1^2-b_0b_2,\ b_2^2-b_0b_1,\ b_0^2-b_1b_2).
\]
The identity
\[
 v_0+v_1+v_2=\frac12\bigl((b_0-b_1)^2+(b_1-b_2)^2+(b_2-b_0)^2\bigr)>0
\]
holds almost surely.  Hence the kernel line meets the positive orthant exactly
when all three coordinates of $v$ are positive.  On the two input orthants in
which the $b_j$ have a common sign this is impossible, because multiplying the
three required strict inequalities gives equality.  On each of the six mixed
sign orthants, two inequalities are automatic and the third is $X^2>YZ$ for
iid half-normal magnitudes.  Therefore
\[
 \mathbb P\{\ker\Phi\text{ meets }\R_{>0}^3\}
 =\frac34\,\mathbb P\{X^2>YZ\}>\frac34\cdot\frac13=\frac14.
\]
The strict inequality follows because
$\{X>\max(Y,Z)\}$ has probability $1/3$, is contained in
$\{X^2>YZ\}$, and the containment is strict on, for example, the positive
probability box
\[
 1.5<X<1.6,\qquad 2<Y<2.1,\qquad 0.5<Z<0.6.
\]
An orthant-invariant random line with no zero coordinates assigns probability
$1/4$ to each of the four projective orthants, so this yields a second
contradiction.  The main theorem is preferable because it lies exactly at
$M=N/2$ and gives a probability-one separator.

\section{Verification, limitations, and novelty}

The accompanying script symbolically expands \eqref{eq:Q}, verifies its
identical vanishing before the sign flip, and verifies that the flipped
polynomial takes value $4$ at the displayed witness.  These checks are not
part of the proof.

The theorem gives a full negative answer to the universal orthant-symmetry
question for the Gaussian partial-circulant ensemble considered by the source,
and it does so in the balanced regime $M=N/2$.  It does not classify every
dimension and row subset, does not settle a fixed-rank Rademacher formulation,
and does not provide the alternative route to the recovery upper bound which
the source requests after a negative answer.

A bounded novelty search on 17 August 2026 covered the four run indexes, the
local parsed source, the current arXiv record, exact searches for the source
question and arXiv id, partial-circulant/orthant-symmetry variants, and the
$2\times3$ cross-product expression.  No later answer or this Pluecker support
obstruction was found.  Novelty confidence is moderate pending specialist
review.

\paragraph{Human review recommendation.}
Check the six minors, the factorization proving $Q=0$, the induced Pluecker
signs under $S$, and the orthogonal-complement transfer.  The proof uses only
elementary algebra and the standard zero-set property of nonzero polynomials.

\begin{thebibliography}{9}
\raggedright
\bibitem{Keiper}
S. Keiper,
\emph{Recovery of Binary Sparse Signals from Structured Biased Measurements},
arXiv:2006.14835 (2020).
\end{thebibliography}

\end{document}
